\documentclass{article}
\usepackage[utf8]{inputenc}
\usepackage[croatian]{babel}
\usepackage[T1]{fontenc}
\usepackage{lmodern}
\usepackage{algorithmic}
\usepackage{algorithm}
\usepackage{longtable}
\usepackage{graphicx}
\usepackage{booktabs}
\usepackage{hyperref}
% Da bi se promjenio stil citiranja umjesto:
% [authoryear, round]
% staviti:
% [numbers, square]
\usepackage[authoryear, round]{natbib}
\usepackage{amsmath}
\usepackage{subfig}
\usepackage{fixltx2e}
\usepackage{todo}
\usepackage{url}
\usepackage{textcomp}

\newcommand{\engl}[1]{(engl.~\emph{#1})}

\begin{document}
\title{Izlučivanje značajki lica Gaborovim filterom}
\author{Tomislav Reicher \and Krešimir Antolić \and Igor Belša \and Marko Ivanković \and Ivan Krišto \and Maja Legac \and Tomislav Novak}

\maketitle

\tableofcontents

\section*{Za raspraviti}
\begin{itemize}
  \item Je li naslov ok? Je li ok spomenuti ``raspoznavanje lica'' ili ``lice''
  u samom naslovu?
\end{itemize}

\section{Uvod}
Raspoznavanje uzoraka je znanstvena disciplina iz područja računarskih znanosti čiji je cilj klasifikacija ili razvrstavanje objekata u jedan od brojnih razreda ili klasa. Iako su područja uporabe brojna u ovom radu koncentacija je na raspoznavanju vizulanih uzoraka, točnije raspoznavanju lica. Prepoznavanje lica ukljućuje računalno prepoznavanje indentiteta na temelju značajki dobivenih obradom slika lica. Iako ljudima lak zadatak, prepoznavanje lica i njihova klasifikacija je veoma zahtjevan posao za računalo. Zadatak postaje tim zahtjevniji ako su lica slikana pod različitim osvijetljenjem, različitim kutovima ili ako osobe na slikama nemaju uvijek isti izraz lica. Kao osnovna motivacija za koriš¡tenje Gaborovog filtera za izvlačenje značajki je veza
sa biološkim osobinama vida kod sisavaca čiji su receptori osjetljivi na
orijentaciju te imaju karakteristične prostorne frekvencije. Gaborov filter može
iskoristit vizualne osobine kao što su lokalizacija prostora, selekcija
orijentacije i karakteristike prostorne frekvencije
\citep{bhuiyan2007onfacerecognition}\nocite{daugman1985uncertainty}.

Model sustava za raspoznavanje uzoraka se može opisati dijagramom
\ref{fig:ru-sustav-model}.

\begin{figure}[htb]
\begin{center}
\begin{picture}(270,130)
\put(20,110){\makebox(40,15){\texttt{Uzorak}}}
\put(58,117){\vector(1,0){15}}
\put(75,110){\framebox(90,15){\textit{Senzor}}}
\put(120,110){\vector(0,-1){15}}
\put(75,80){\framebox(90,15){\textit{Pretprocesiranje}}}
\put(120,80){\vector(0,-1){15}}
\put(75,50){\framebox(90,15){\textit{Izlučivanje značajki}}}
\put(120,50){\vector(0,-1){15}}
\put(75,20){\framebox(90,15){\textit{Pravilo odlučivanja}}}
\put(167,27){\vector(1,0){15}}
\put(180,20){\makebox(40,15){\texttt{Razred}}}
\end{picture}
\caption{Model sustava za raspoznavanje uzoraka}
\label{fig:ru-sustav-model}
\end{center}
\end{figure}

% Ovo je opis o čemu koji odjeljak govori. Ovo se MORA nalaziti na kraju uvoda. 
U 2.~odjeljku prikazan je matematički model dvodimenzionalnog gaborovog filtera,
u 3.~odjeljku objašnjen je način izvlačenja značajki, 4.~odjeljak navodi objašnjenja
pojedinih parametara i njihov utjecaj na krajnji rezultat, a 5.~prikazuje
rezultate primjene gaborovog filtera na neke od uzoraka. U 6.~odjeljku je
objašnjena interpretacija značajki te način korištenja dobivenih značajki u
sustavima za raspoznavanje uzoraka, odnosno klasifikacija pomoću izlučenih
značajki. Zaključak je dan u 7.~odjeljku.\TODO{Prešturo\ldots Doraditi nakon što
budemo znali što ćemo uopće pisati\ldots}

\section{Dizajn Gaborovog filtera}
Dvodimenzionalna Gaborova funkcija data je kao \citep{petkovgabor}:
\begin{equation}
% g(x,y)=\frac{1}{2\pi \sigma^2_{xy}}e^{-\left ( \frac{x'^2 +
% y'^2}{2\sigma^2_{x,y}} \right)} \times \left ( e^{2\pi i r_0 x'} -
% e^{-\frac{r_0^2}{2\sigma^2_{uv}}}\right),
g_{\lambda,\theta,\varphi,\sigma,\gamma}(x,y) = \exp\left ( -
\frac{x'^2+\gamma^2 y'^2}{2\sigma^2}\right ) \cos \left ( 2\pi \frac{x'}{\lambda} + \varphi \right ),
\label{2d-gabor}
\end{equation}
pri čemu su
\begin{eqnarray*}
x' = x \cos \theta + y \sin \theta, \\
y' = -x \sin \theta + y \cos \theta.
\end{eqnarray*}
% gdje je $\sigma_{xy}$ standardna devijacija Gaussove omotnice\TODO{\emph{Gaussian
% envelope}?}koja karakterizira prostorni obujam i širinu filtera. 
Frekvencija i odabir orijentacije Gaborovog filtera su izražajnije u domeni
frekvencijskog prikaza predstavljenog jednadžbom (\ref{gabor-frek}) koja
određuje koliko filter utječe na svaku frekvencijsku komponentu ulazne slike.
\begin{equation}
G(u,v) = \exp \left ( - \frac{(u-u_0)^2 + (v-v_0)^2}{2\sigma^2_{uv}}\right ) -
\exp \left ( - \frac{r_0^2}{2\sigma^2_{uv}} \right),
\label{gabor-frek}
\end{equation}
\begin{equation}
\sigma_{uv} = \frac{1}{2\pi \sigma_{xy}}.
\end{equation}
Parametri ($u_0$, $v_0$) definiraju prostornu frekvenciju sinusoidalnog vala u
ravnini koji također može biti izražen polarnim koordinatama kao radialna frekvencija $r_0$ i
orijentacija $\theta$:
\begin{eqnarray}
r_0^2 = u_0^2 + v_0^2, \\
\tan \theta = \frac{v_0}{u_0}.
\end{eqnarray}
Osobina Gaborovog filtera definirana je radijalnom frekvencijom $r_0$,
orijentacijom i širinom filtera.

Ako je svrha filtera izvlačenje značajki lica, zanimaju nas četiri orijentacije
\citep{huang2005robust}:
\begin{equation}
\theta_k = \frac{\pi(k-1)}{4},\: k = 1,2,3,4.
\label{equ:orijentacije}
\end{equation}
Orijentacije se odnose na oblu konturu lica, oči i usta koji se nalaze u skoro
horizontalnoj ravnini, te nos koji je u vertikalnoj ravnini.

\section{Izvlačenje značajki pomoću gaborovog filtera}
Gaborove značajke se dobivaju konvolucijom klizečeg prozora slike i Gaborovog
filtera.

Neka je $I(x,y)$ slika. Konvolucija slike $I(x,y)$ i Gaborovog filtera dana je
sa:
\begin{equation}
O(x,y,r_0, \theta_k) = I(x,y) * g(x,y,r_0, \theta_k),
\label{konvolucija-filter-slika}
\end{equation}
pri čemu je $k = 1, 2, 3, 4$. $O(x,y,r_0, \theta_k)$ nazivamo Gaborov prikaz
slike $I(x,y)$.\TODO{Treba nastaviti\ldots pogledati \citep{huang2005robust}
dio 3.2.}

\subsection{Preprocesiranje uzoraka}
\TODO{Izvuči iz \citep{bhuiyan2007onfacerecognition}}

\section{Utjecaj pojedinih parametara}
Gaborov filter definiraju valna duljina ($\lambda$),
orijentacija ($\theta$), fazni pomak ($\varphi$), omjer dimenzija \engl{aspect ratio}
($\gamma $) i širina filtera \engl{bandwidth} ($b$) \citep{petkovgabor}.
Promjena parametara jest stvaranje novog filtera.\TODO{Ovo provjeriti da nisam pokupio
nešto karakteristično isključivo za onaj program iz onog HTML--a čiji je link
naveden dolje u literaturi}

\subsection{Valna duljina ($\lambda$)}
Valna duljina se odnosi na valnu duljinu kosinusa u Gaborovoj funkciji.
Vrijednost je realan broj veći ili jednak $2$. Jedinica su pixeli.

Slike Gaborovih filtera uz parametre $\varphi, \theta = 0,\: b = 1$, i $\gamma =
0.5$ možete vidjeti na slici \ref{fig:filter-wavelengths}.

\begin{figure}[htb]
  \centering
  \subfloat[Filter uz $\lambda = 5$.]{
  	\label{fig:filter-lambda-5}
  	\includegraphics[width=4cm]{images/wavelength5.jpg}
  }
  \hspace{50pt}
  \subfloat[Filter uz $\lambda = 10$.]{
  	\label{fig:filter-lambda-10}
  	\includegraphics[width=4cm]{images/wavelength10.jpg}
  }
  \caption{Gaborov filter uz različite valne duljine.}
  \label{fig:filter-wavelengths}
\end{figure}

% TODO: Pogledati o čemu se ovdje radi\ldots Gledao sam formulu, ništa pametno
% nisam zaključio, generirao sam slike, također ništa pametno\ldots Nešto mi
% promiče\ldots Zakomentirao sam tu rečenicu dok mi ne postane smislena.
%Uz $\lambda = 2$, ne bi se smjela koristiti kombinacija $\varphi = \pm 90$ jer
%u tom slučaju Gaborova funkcija je uzrokovana u svojim
%nul--prijelazima\TODO{``sampled in its zero crossings''?}.
Da bi se spriječila pojava neželjenih efekata na rubovima slike, vrijednost valne
duljine mora biti manja od $\frac{1}{5}$ veličine slike koja prolazi kroz
filter. Filter sa prevelikom valnom duljinom možete vidjeti na slici
\ref{fig:wavelength-overflow}. U datom primjeru filter djeluje na području
širine 100 piksela, a filter je širok 30 piksela.

\begin{figure}[htb]
\begin{center}
\includegraphics[width=4cm]{images/wavelength30overflow.jpg}
\end{center}
\caption{Prikaz filtera uz preveliku valnu duljinu.}
\label{fig:wavelength-overflow}
\end{figure}

\subsection{Orijentacija ($\theta$)}
Orijentacija određuje kut između normale i paralelnih pruga Gaborovog
filtera.\TODO{Što je normala, a što su te pruge\ldots nisam uspio skužiti\ldots}
Određena je kutom od $0$ do $360$ stupnjeva. Prikaze Gaborovog filtera sa
orijentacijama od 45\textdegree, 80\textdegree\,i 0\textdegree\, možete vidjeti
na slikama \ref{fig:filter-orientation-45}, \ref{fig:filter-orientation-80} i
\ref{fig:filter-lambda-10}.

\begin{figure}[htb]
  \centering
  \subfloat[Filter uz $\theta = 45$\textdegree.]{
  	\label{fig:filter-orientation-45}
  	\includegraphics[width=4cm]{images/orientation45.jpg}
  }
  \hspace{50pt}
  \subfloat[Filter uz $\theta = 80$\textdegree.]{
  	\label{fig:filter-orientation-80}
  	\includegraphics[width=4cm]{images/orientation80.jpg}
  }
  \caption{Gaborov filter uz različite orijentacije.}
  \label{fig:filter-orientations}
\end{figure}

\subsection{Fazni pomak ($\varphi$)}
Ovaj parametar određuje fazni pomak funkcije kosinus unutar Gaborove funkcije.
Ispravne vrijednosti su realni brojevi između $-180$ i $180$ stupnjeva. Prikaze
Gaborovog filtera sa fazama od 45\textdegree, 180\textdegree\, i 0\textdegree\,
možete vidjeti na slikama \ref{fig:filter-phase-45},
\ref{fig:filter-phase-180} i \ref{fig:filter-lambda-10}.

\begin{figure}[htb]
  \centering
  \subfloat[Filter uz $\varphi = 45$\textdegree.]{
  	\label{fig:filter-phase-45}
  	\includegraphics[width=4cm]{images/phase45.jpg}
  }
  \hspace{50pt}
  \subfloat[Filter uz $\varphi = 180$\textdegree.]{
  	\label{fig:filter-phase-180}
  	\includegraphics[width=4cm]{images/phase180.jpg}
  }
  \caption{Gaborov filter uz različite faze.}
  \label{fig:filter-phases}
\end{figure}

Vrijednosti $0$\textdegree\, i $180$\textdegree\, odgovaraju
središnje--simetričnim funkcijama, a $-90$\textdegree\, i $90$\textdegree\,
anti--simetričnim funkcijama (primjer filtera se može vidjeti na slici
\ref{fig:phase-90}). Svi ostali slučajevi odgovaraju asimetričnim funkcijama.

\begin{figure}[htb]
\begin{center}
\includegraphics[width=4cm]{images/phase90.jpg}
\end{center}
\caption{Filter uz fazu od 90\textdegree (anti--simetrija).}
\label{fig:phase-90}
\end{figure}

\subsection{Omjer dimenzija ($\gamma $)}
Parametar, koji se preciznije naziva prostorni omjer dimenzija, koji određuje
eliptičnost Gaborove funkcije. Za $\gamma = 1$, eliptičnost se svodi na krug.
Za $\gamma < 1$ funkcija je izdužena u smjeru paralelnom s paralelnim prugama
funkcije\TODO{Opet te paralelne pruge!}. Primjer filtera sa različitim omjerima
dimenzija se može vidjeti na slici \ref{fig:filter-ratios}.

\begin{figure}[htb]
  \centering
  \subfloat[Filter uz $\gamma = 0.5$.]{
  	\label{fig:filter-gamma-05}
  	\includegraphics[width=4cm]{images/wavelength10.jpg}
  }
  \hspace{50pt}
  \subfloat[Filter uz $\gamma = 1$.]{
  	\label{fig:filter-gamma-1}
  	\includegraphics[width=4cm]{images/ratio1.jpg}
  }
  \caption{Gaborov filter uz različite omjere dimenzija.}
  \label{fig:filter-ratios}
\end{figure}

\subsection{Širina filtera ($b$)}
Prostorna širina filtera $b$ (u oktavama) Gaborovog filtera je povezana sa
omjerom $\frac{\sigma}{\lambda}$, pri čemu su $\sigma$ i $\lambda$ standardna
devijacija Gaussovog umnoška Gaborove funkcije i valna duljina. Definirana je
sa:
\begin{eqnarray}
b = \log_2{\left (  \frac{\frac{\sigma}{\lambda}\pi + \sqrt{\frac{\ln2}{2}}}
{\frac{\sigma}{\lambda}\pi - \sqrt{\frac{\ln2}{2}}} \right )}, \\
\frac{\sigma}{\lambda} =
\frac{1}{\pi}\sqrt{\frac{ln2}{2}}\cdot\frac{2^b+1}{2^b-1}.
\end{eqnarray}
Vrijednost $\sigma$ se ne može direktno odrediti. Ona se može mijenjati samo
preko vrijednosti širine filtera, $b$.

Širina filtera se određuje kao realni pozitivni broj. Što je širina manja,
$\sigma$ je veća i povećava se broj naglašavajućih i prigušavajućih pruga
Gaborovog filtera.\TODO{Je li ovo ok ovako napisano?} Primjere filtera sa
različitim širinama možete vidjeti na slikama \ref{fig:filter-bandwidth-05} i
\ref{fig:filter-bandwidth-2}. Na prethodnim primjerima korištena je širina
filtera od $b = 1$.

\begin{figure}[htb]
  \centering
  \subfloat[Filter uz $b = 0.5$.]{
  	\label{fig:filter-bandwidth-05}
  	\includegraphics[width=4cm]{images/bandwidth05.jpg}
  }
  \hspace{50pt}
  \subfloat[Filter uz $b = 2$.]{
  	\label{fig:filter-bandwidth-2}
  	\includegraphics[width=4cm]{images/bandwidth2.jpg}
  }
  \caption{Gaborov filter uz različite širine.}
  \label{fig:filter-bandwidths}
\end{figure}

\section{Primjer primjene}
Izvuči iz http://matlabserver.cs.rug.nl/edgedetectionweb/web/edgedetection\_examples.html

\section{Klasifikacija}

\section{Zaključak}

\bibliography{literatura}
\bibliographystyle{plainnat}

\newpage
\appendix
\section{\LaTeX~playground}
Ovdje se možete igrati sa \LaTeX--om. Ideja je da igranjem u ovome dijelu
naučite nešto korisno oko \LaTeX--a tako da vam oni koji znaju \LaTeX~tu ostave
koji koristan primjer uporabe. Uglavnom, uvijek se igrajte s \LaTeX--om, jer igranje bez \LaTeX--a nije sigurno igranje. 

\subsection{O novim redovima}
Dokument dijelimo na paragrafe. Tekst unutar paragrafa se ne razlama. Ilitiga,
ako napišete:
\begin{verbatim}
Želim ovo u prvom,
ovo u drugom,
a ovo u trećem, i da je blizu udaljena od       daleko.
\end{verbatim}
Dobiti će te:
\emph{Želim ovo u prvom,
ovo u drugom,
a ovo u trećem, i da je blizu udaljena od       daleko.}

Ako želite novi red unutar paragrafa, morat će te dodati
\verb|\\| na kraj red ili komandu \verb|\newline|. I to je ružno\ldots

Ako želite preći u novi paragraf dovoljno je ostaviti jedan prazan red između
prošlog i novog paragrafa, tj.:
\begin{verbatim}
Prvi paragraf.

Drugi paragraf.
\end{verbatim}

\subsection{Korištenje TODO i FUJ naredbi}
Za potrebe ovog rada, dodane su \verb|\TODO{}| i \verb|\FUJ{}| naredbe. Služe
da bi popljuvali\FUJ{Kako ružna riječ}nečije djelo ili napisali što još treba
napraviti\TODO{Navedi primjer korištenja! :D}.

U argumente ovih naredbi, tj. sadržaj \textsf{TODO}--a možete ugurati sam kod.
Npr.~recimo.\TODO{\LaTeX, ovo je \emph{naglašeno}; ili
malo matematike $\iint_a^b{x^2dx}$} No, ovdje neće proći stvari kao više
linijski programski kod i sl.

\subsection{O pisanju matematičkih izraza}
Ako želite naučiti koristiti \LaTeX~za ono što ovaj podnaslov spominje, trebati
će vam neki tutorial ili knjiga (a ima ih masu\ldots i to besplatnih!).
Ja ću samo navesti par osnova.

Primjer izraza koji se nalazi u posebnom bloku, i još je centriran:
$$x_{1,2} = \frac{-b \pm \sqrt{b^2-4ac}}{2a}$$
Napisao sam: \verb|$$x_{1,2} = \frac{-b \pm \sqrt{b^2-4ac}}{2a}$$|

Ovo je jedan način zapisa. Mogli smo koristiti:
\begin{equation}
x_{1,2} = \frac{-b \pm \sqrt{b^2-4ac}}{2a}
\label{rjesenje-kvadratne-jed}
\end{equation}
Napisao sam:
\begin{verbatim}
\begin{equation}
x_{1,2} = \frac{-b \pm \sqrt{b^2-4ac}}{2a}
\label{rjesenje-kvadratne-jed}
\end{equation}
\end{verbatim}

Na ovaj drugi način izrazi se numeriraju i možemo se lako referencirati na
njih, npr.~(\ref{rjesenje-kvadratne-jed}). Referenciramo se ne bilo koji
\verb|\label{}| pomoću naredbe \verb|\ref{}|,
npr.~\verb|\ref{rjesenje-kvadratne-jed}|.

Evo neka stranica sa par zanimljivih primjera koje najvjerojatnije nikad nećete
imati priliku primjeniti, \url{http://www.personal.ceu.hu/tex/cookbook.html} te
još malo filozofije o svemu tome
\url{http://www.math.uiuc.edu/~hildebr/tex/displays.html}.

Btw.~ako gledate source, možda se pitate zašto uvijek stavljam ``\textasciitilde{}''
nakon točke. Radi se o tome da \LaTeX~nakon svake točke stavi malo više
razmaka, jer misli da se radi o početku nove rečenice. Ako mu stavimo tildu,
onda taj razmak bude onakav kakav bi trebao biti.

\subsection{O literaturi}
Literaturu izvlačite sa citeseerxa ili google schoolara jer vam oni odmah daju
i bibtex članka koji samo kopirate u file \emph{literatura.bib}.

Primjer citiranja: \verb|\citep{Yang04facerecognition}|.

\end{document}
